\chapter{先建立执行模型}

\section{从一个错误的猜测开始}

很多入门材料会直接讲 \code{int}、\code{if}、\code{for}、\code{class}。这些概念当然重要，但如果一开始不回答“程序运行时到底发生了什么”，后面的指针、引用、移动语义和 RAII 都会像孤立规则。我们先看一个最小程序：

\begin{lstlisting}[language=C++,caption={一个看似普通的程序}]
#include <iostream>
#include <string>

int main()
{
    std::string name = "Ada";
    int score = 95;
    std::cout << name << " " << score << '\n';
    return 0;
}
\end{lstlisting}

这段程序能输出 \code{Ada 95}。如果只从语法角度看，它包含头文件、变量、输出和返回值。但从执行模型看，它还包含更多信息：\code{name} 是一个对象，它管理一段字符资源；\code{score} 是一个整数值；\code{std::cout} 是标准库提供的输出流；程序从 \code{main} 开始执行，函数结束时局部对象按逆序销毁。

\begin{mentalmodel}
\cpp 程序可以先用四个问题观察：对象在哪里存在，谁拥有资源，什么时候构造和析构，错误如何被发现和传播。语法只是表达这些事实的工具。
\end{mentalmodel}

\section{编译不是运行}

初学者常把“能编译”理解成“程序正确”。这不够。编译器主要检查语法、类型和一部分静态规则；运行时才会遇到输入、文件、内存、时间和并发调度。下面的程序能编译，但它不可靠：

\begin{lstlisting}[language=C++,caption={能编译但会越界的程序}]
#include <iostream>
#include <vector>

int main()
{
    std::vector<int> scores{90, 85, 77};
    std::cout << scores[3] << '\n';
}
\end{lstlisting}

\code{scores[3]} 访问第四个元素，但数组只有三个元素。标准库的 \code{operator[]} 通常不会做边界检查，程序可能输出杂乱值，也可能崩溃。改成 \code{scores.at(3)} 后，标准库会在越界时抛出异常。这里不是说所有代码都必须用 \code{at}，而是要理解：接口选择影响错误出现的位置。

\begin{warningbox}
能编译只说明程序通过了语言规则检查，不说明数据范围、资源生命周期和业务不变量正确。写 \cpp 时必须同时关心编译期和运行期。
\end{warningbox}

\section{从源码到可执行文件}

一个常见构建链条是：预处理、编译、汇编、链接。预处理处理 \code{#include} 和宏；编译把翻译单元转换成目标代码；链接把多个目标文件和库合成可执行文件。先记住这个模型，后面遇到“头文件重复定义”“链接找不到符号”“模板为什么要放头文件里”时就不会只靠猜。

\begin{lstlisting}[language=bash,caption={手动观察构建步骤}]
g++ -std=c++20 -E main.cpp -o main.i
g++ -std=c++20 -S main.cpp -o main.s
g++ -std=c++20 -c main.cpp -o main.o
g++ main.o -o app
\end{lstlisting}

第一步生成的 \code{main.i} 会很大，因为头文件内容被展开进来。第三步生成 \code{main.o}，它还不是最终程序。第四步链接时，链接器需要找到 \code{std::cout}、字符串构造函数等符号的实现。

\section{对象、值和名字}

变量名不是对象本身，而是访问对象的名字。下面的代码里，\code{x} 和 \code{y} 是两个不同对象：

\begin{lstlisting}[language=C++,caption={复制会产生新对象}]
int x = 10;
int y = x;
y = 20;
\end{lstlisting}

执行后 \code{x} 仍然是 \code{10}，\code{y} 变成 \code{20}。如果换成引用，情况就不同：

\begin{lstlisting}[language=C++,caption={引用是已有对象的别名}]
int x = 10;
int& alias = x;
alias = 20;
\end{lstlisting}

这时 \code{x} 也变成 \code{20}。学习 \cpp 的第一个分水岭就在这里：你必须看清一行代码是在创建新对象、观察旧对象，还是通过别名修改旧对象。

\begin{keyidea}
名字、对象和值要分开理解。名字是访问入口，对象有生命周期，值是对象当前保存的信息。引用不创建新对象，它让一个已有对象有了另一个访问入口。
\end{keyidea}

\section{本章边界检查}

读完本章后，你应该能回答这些问题：

\begin{itemize}
  \item \code{std::string name = "Ada"} 创建了什么对象？
  \item 为什么 \code{scores[3]} 能编译却不可靠？
  \item 编译错误和链接错误有什么区别？
  \item 复制和引用在对象层面有什么不同？
\end{itemize}

后面的章节会不断回到这些问题。只要执行模型稳，语法规则就不再是散点。

\section{把一行程序拆成执行步骤}

为了让这个模型不只停留在口号上，我们把第一段程序里的两行拆开看：

\begin{lstlisting}[language=C++,caption={两行代码里的对象变化}]
std::string name = "Ada";
std::cout << name << '\n';
\end{lstlisting}

第一行不是“变量等于字符串”这么简单。更准确的过程是：编译器看到左边类型是 \code{std::string}，右边是字符串字面量；于是构造一个 \code{std::string} 对象，让这个对象保存字符内容，并把名字 \code{name} 绑定到这个对象。第二行也不是把对象复制到屏幕上，而是调用输出流的插入操作，把 \code{name} 当前表示的字符序列写到标准输出。

如果把 \code{name} 改成整数，输出语句仍然长得差不多：

\begin{lstlisting}[language=C++,caption={同一个输出符号背后是不同操作}]
int score = 95;
std::cout << score << '\n';
\end{lstlisting}

这说明 \code{<<} 不是一个固定的机器动作，而是一组根据类型选择的操作。初学时不用急着理解“重载”的全部规则，但要先记住：\cpp 代码表面相似，不代表执行细节相同。类型决定能做什么，也决定怎么做。

\section{三类错误要分开定位}

写程序时常见三类错误。第一类是编译错误，例如少分号、类型不匹配、调用不存在的函数。第二类是链接错误，例如声明了函数却没有提供实现。第三类是运行错误，例如越界、除零、文件不存在、输入格式不对。

\begin{lstlisting}[language=C++,caption={声明了函数但没有实现会导致链接错误}]
int add(int lhs, int rhs);

int main()
{
    return add(1, 2);
}
\end{lstlisting}

这段代码能通过语法检查，因为编译器知道 \code{add} 的参数和返回类型。但链接时找不到函数体，所以会失败。许多初学者看到一屏英文报错就慌，其实只要先分类，定位就清楚很多：编译错看当前源文件和类型，链接错看函数实现和目标文件，运行错看输入、数据范围和生命周期。

\section{一个最小实验流程}

建议从第一章开始就固定一个实验流程。把下面代码保存为 \filepath{main.cpp}：

\begin{lstlisting}[language=C++,caption={最小实验程序}]
#include <iostream>

int main()
{
    int value = 41;
    ++value;
    std::cout << value << '\n';
}
\end{lstlisting}

用下面命令构建并运行：

\begin{lstlisting}[language=bash,caption={手动编译运行}]
g++ -std=c++20 -Wall -Wextra -pedantic main.cpp -o app
./app
\end{lstlisting}

\code{-Wall -Wextra -pedantic} 会打开更多警告。警告不是噪声，它经常指出“能编译但值得怀疑”的代码。以后每写一个小程序，都先做到三件事：能用固定命令构建，能说明输出为什么是这个值，能解释如果输入为空或越界会怎样。

\begin{exercisebox}
把 \code{scores[3]} 改成 \code{scores.at(3)}，重新运行并观察错误出现在哪里。再把 \code{std::string name = "Ada"} 改成 \code{auto name = "Ada"}，思考这时 \code{name} 的类型还是不是 \code{std::string}。这两个小实验分别训练运行期边界和类型推导。
\end{exercisebox}
